\clearpage
\chapter*{Errata da exposição n.\textsuperscript{o} 195}
\addcontentsline{toc}{chapter}{Errata da exposição n.\textsuperscript{o} 195}
\setcounter{footnote}{0}
\src{300}

\noindent\textbf{4. TDTE II: O teorema de existência na teoria formal dos
módulos.}

\begin{center}
[Seminário Bourbaki, t. 12, 1959/60, n.\textsuperscript{o} 195, 22 p.]
\end{center}

\noindent\textbf{Página 8.} --- A fórmula escrita na proposição 5.1 só é
correta quando $\Lambda$ é um corpo; no caso geral, deve-se substituir
$\mathfrak m_\xi/\mathfrak m_\xi^2$ pelo quociente deste espaço pela
imagem de $\mathfrak n_\xi/\mathfrak n_\xi^2$, onde $\mathfrak n$ é o ideal
maximal de $\Lambda$. Além disso, a definição dada para que $O$ seja liso sobre
$\Lambda$ só é correta quando a extensão residual $k'/k$ é separável.
No caso geral, cf. SGA III, 1.1.

\noindent\textbf{Página 14, observação.} --- Os problemas formulados aqui
ficam completamente resolvidos no caso projetivo mediante os «esquemas de
Hilbert» (cf. TDTE IV). Exemplos de Nagata e Hironaka mostram, em contrapartida, que
os funtores considerados não são necessariamente representáveis se não se impõe
uma hipótese de projetividade, mesmo limitando-se à classificação das
subvariedades de dimensão zero de uma variedade completa não singular de
dimensão 3; isso está ligado ao fato de que o quadrado simétrico de uma
tal variedade pode não existir.

\src{301}
\noindent\textbf{Página 15, n.\textsuperscript{o} 3.} --- Para um estudo mais
completo, veja-se TDTE V.

\noindent\textbf{Página 16, proposição 3.1.} --- Em vez de «se $f$ é
próprio», leia-se «se $f$ é próprio e separável».

\noindent\textbf{Páginas 18 e 19, observações, 2.\textsuperscript{o}} --- Demonstrei
recentemente que o esquema formal de módulos de uma variedade
abeliana sobre um corpo é efetivamente liso sobre o anel de Witt; em
outros termos, todo esquema abeliano sobre um anel artiniano local que seja
quociente de outro provém, por redução, de um esquema abeliano sobre este
último. A demonstração utiliza unicamente as propriedades de variância da
classe de obstrução ao levantamento introduzida na exposição
n.\textsuperscript{o} 182, página 12. Recordemos igualmente que os esquemas de
módulos das curvas de gênero $g$ e dos esquemas abelianos polarizados foram
construídos por Mumford (cf. SHMT).

\noindent\textbf{Página 19, parágrafo 5.} --- Deve-se supor que a seção
considerada para definir $U$ seja um não divisor de zero não somente sobre
$\mathfrak X$, mas também sobre todo $X_n$.

\noindent\textbf{Página 20, linhas 12 e 13.} --- Em vez de
\[
  G_x=G\bigl(\operatorname{Spec}(\widehat{\mathcal O}_{\mathfrak X,x})\bigr),
\]
leia-se
\[
  G_x=G\bigl(\operatorname{Spec}(\mathcal O_{\mathfrak X,x})\bigr)
  \subset
  G\bigl(\operatorname{Spec}(\widehat{\mathcal O}_{\mathfrak X,x})\bigr),
\]
e, em vez de «$\mathcal O'_{0x}$», leia-se
«$\widehat{\mathcal O}'_{0x}$».
